传统三维机械建模软件通常要求用户具备较强的参数化建模经验，这使教学演示、课程设计和简单工装场景中的建模过程仍然高度依赖人工操作，并造成从文本需求到三维表达之间的转换成本偏高。现有自然语言驱动建模方案虽然显示出应用潜力，但若直接依赖大语言模型自由生成 CAD 脚本，往往难以同时保证语法正确性、几何语义一致性和工程可验证性。针对这一问题，本文设计并实现了面向简单机械结构建模的自然语言驱动参数化 CAD 生成系统 \texttt{fs-builder}，采用“需求分析、结构化 Plan 校验、确定性 FeatureScript 模板生成”的三级处理链路，并提供 CLI 与本地 Web UI 两类交互入口。基于项目当前实现，系统支持 \texttt{box}、\texttt{cylinder}、\texttt{hollow\_cylinder}、\texttt{tapered\_cylinder}、\texttt{flange} 五种基础零件形状；自动化测试结果显示 50 项测试全部通过，总代码覆盖率达到 82\%，在冷拉延模具样例上系统成功识别 5 个主件与 4 条装配关系，并完成 5/5 零件的脚本生成。本文的独特贡献在于提出了一种适用于教学与原型验证场景的工程化路线，即以大模型负责需求理解、以强类型中间表示承担约束收缩、以确定性模板保障最终 CAD 脚本可控，从而在自然语言灵活性与参数化建模稳定性之间建立可实现的平衡。
